%% preambolo.tex
%%
%% Copyright (C) 2000-2018 Simone Piccardi.  Permission is granted to
%% copy, distribute and/or modify this document under the terms of the GNU Free
%% Documentation License, Version 1.1 or any later version published by the
%% Free Software Foundation; with the Invariant Sections being "Un preambolo",
%% with no Front-Cover Texts, and with no Back-Cover Texts.  A copy of the
%% license is included in the section entitled "GNU Free Documentation
%% License".
%%

\chapter{Un preambolo}
\label{cha:preamble}

Questa guida nasce dalla mia profonda convinzione che le istanze di libertà e
di condivisione della conoscenza che hanno dato vita a quello straordinario
movimento di persone ed intelligenza che va sotto il nome di \textsl{software
  libero} hanno la stessa rilevanza anche quando applicate alla produzione
culturale in genere.

L'ambito più comune in cui questa filosofia viene applicata è quello della
documentazione perché il software, per quanto possa essere libero, se non
accompagnato da una buona documentazione che aiuti a comprenderne il
funzionamento, rischia di essere fortemente deficitario riguardo ad una delle
libertà fondamentali, quella di essere studiato e migliorato.

Ritengo inoltre che in campo tecnico ed educativo sia importante poter
disporre di testi didattici (come manuali, enciclopedie, dizionari, ecc.) in
grado di crescere, essere adattati alle diverse esigenze, modificati e
ampliati, o anche ridotti per usi specifici, nello stesso modo in cui si fa
per il software libero.

Questa guida è il mio tentativo di restituire indietro, nei limiti di quelle
che sono le mie capacità, un po' della conoscenza che ho ricevuto, mettendo a
disposizione un testo che possa fare da riferimento a chi si avvicina alla
programmazione su Linux, nella speranza anche di trasmettergli non solo delle
conoscenze tecniche, ma anche un po' di quella passione per la libertà e la
condivisione della conoscenza che sono la ricchezza maggiore che ho ricevuto.

E, come per il software libero, anche in questo caso è importante la
possibilità di accedere ai sorgenti (e non solo al risultato finale, sia
questo una stampa o un file formattato) e la libertà di modificarli per
apportarvi migliorie, aggiornamenti, ecc.

Per questo motivo la Free Software Foundation ha creato una apposita licenza
che potesse giocare lo stesso ruolo fondamentale che la GPL ha avuto per il
software libero nel garantire la permanenza delle libertà date, ma potesse
anche tenere conto delle differenze che comunque ci sono fra un testo ed un
programma.

Una di queste differenze è che in un testo, come in questa sezione, possono
venire espresse quelle che sono le idee ed i punti di vista dell'autore, e
mentre trovo che sia necessario permettere cambiamenti nei contenuti tecnici,
che devono essere aggiornati e corretti, non vale lo stesso per l'espressione
delle mie idee contenuta in questa sezione, che ho richiesto resti invariata.

Il progetto pertanto prevede il rilascio della guida con licenza GNU FDL, ed
una modalità di realizzazione aperta che permetta di accogliere i contributi
di chiunque sia interessato. Tutti i programmi di esempio sono rilasciati con
licenza GNU GPL.




%%% Local Variables: 
%%% mode: latex
%%% TeX-master: "gapil"
%%% End: 

% LocalWords:  Foundation FDL
